\begin{frame}[plain]
  \centering
  \vspace*{0.18\paperheight}
  \LARGE \textbf{Fine-Tuning Methods for LLMs}
\end{frame}

\begin{frame}{Categories of Fine-Tuning}
  \begin{center}
    \begin{tabular}{|l|p{3.3cm}|c|c|}
      \hline
      \textbf{Method} & \textbf{Description} & \textbf{Params} & \textbf{Efficiency} \\
      \hline
      Full FT & Retrain all weights & All & Expensive \\
      Adapter & Add small bottlenecks (Houlsby) & Few & Efficient \\
      Prefix & Tune soft prompts & Few tokens & Efficient \\
      LoRA/QLoRA & Low-rank weight deltas & Very few & Very Efficient \\
      Instr. Tuning & Tune on instr. datasets & All/Part & - \\
      \hline
    \end{tabular}
  \end{center}
\end{frame}

\begin{frame}[t]{Full Fine-Tuning (FT)}
  \begin{itemize}
    \item[\textcolor{red}{$\blacktriangleright$}] \textbf{Definition:} Fine-tune all parameters of the pre-trained model on downstream data.
    \item[\textcolor{red}{$\blacktriangleright$}] \textbf{Historical Use:} Common in early GPT-2 and BERT applications.
    \item[\textcolor{red}{$\blacktriangleright$}] \textbf{Pros:}
      \begin{itemize}
        \item Maximum flexibility and performance
      \end{itemize}
    \item[\textcolor{red}{$\blacktriangleright$}] \textbf{Cons:}
      \begin{itemize}
        \item Expensive (requires large compute resources)
        \item Prone to catastrophic forgetting
        \item Not efficient for large models
      \end{itemize}
  \end{itemize}
\end{frame}

\begin{frame}{Full Fine-Tuning Workflow}
  \includegraphics[width=0.60\paperwidth,height=0.60\paperheight,keepaspectratio]{assets/Lecture-8_Fine_Tuning_RLHF/slide_010.png}
\end{frame}

\begin{frame}{Parameter-Efficient Fine-Tuning (PEFT)}
  \begin{itemize}
    \item[\textcolor{red}{$\blacktriangleright$}] \textbf{Key Idea:} Keep the base model frozen, tune only a small subset of parameters.
    \item[\textcolor{red}{$\blacktriangleright$}] \textbf{Types:}
    \begin{itemize}
      \item Adapter Modules (Houlsby et al., 2019)
      \item Prompt Tuning / Prefix Tuning
      \item LoRA / QLoRA
    \end{itemize}
    \item[\textcolor{red}{$\blacktriangleright$}] \textbf{Benefits:}
    \begin{itemize}
      \item Enables multi-tasking and personalization
      \item Suitable for low-resource adaptation
      \item Reduces compute and memory requirements
    \end{itemize}
  \end{itemize}
\end{frame}


\begin{frame}{Parameter-Efficient Fine-Tuning (PEFT)}
  \centering
  \includegraphics[width=0.60\paperwidth,height=0.60\paperheight,keepaspectratio]{assets/Lecture-8_Fine_Tuning_RLHF/slide_012.png}
\end{frame}
